%% kapa: exam=1 round=29 year=2018 period=2 form=A subject=accounting subject_ko=회계학 seq=3 number=43
%% source: Q-Net 감정평가사 기출문제 (cst00309) artlSeq=5202808 file=extracted/1008/5202808/2018년 제29회 감정평가사 1차시험 기출문제/2018년 제29회 감정평가사 1차시험 기출문제/2교시 A형.pdf
%% conversion: kordoc 4.24.51 -> tools/gen1.py (auto); review=proofread-sonnet; flags=none
\begin{question}{43}
(주)감평은 20x1년 초 기계장치를 ￦100,000에 취득하고 재평가모형을 적용하기로 하였다. 기계장치의 내용연수는 5년, 잔존가치는 ￦0이며 정액법으로 감가상각 한다. 다음은 연도별 기계장치의 공정가치와 회수가능액이다.
\par{\small\centering
\begin{tabular}{lrr}
\hline
 & 20x1년 말 & 20x2년 말 \\
공정가치 & ￦88,000 & ￦60,000 \\
회수가능액 & 90,000 & 48,000 \\
\hline
\end{tabular}\par}
(주)감평은 20x2년 말에 기계장치에 대해서 손상이 발생하였다고 판단하였다. 기계장치와 관련하여 20x2년도 포괄손익계산서 상 당기순이익과 기타포괄이익에 미치는 영향은 각각 얼마인가? (단, 기계장치를 사용함에 따라 재평가잉여금의 일부를 이익 잉여금으로 대체하지 않는다.) \par
\hspace{2.5em}\makebox[6cm][l]{\uline{당기순이익에 미치는 영향}}\uline{기타포괄이익에 미치는 영향} \par
\begin{choices}
  \item \makebox[6cm][l]{￦10,000 감소}￦2,000 증가
  \item \makebox[6cm][l]{￦10,000 감소}￦8,000 감소
  \item \makebox[6cm][l]{￦32,000 감소}￦2,000 감소
  \item \makebox[6cm][l]{￦32,000 감소}￦2,000 증가
  \item \makebox[6cm][l]{￦32,000 감소}￦8,000 감소
\end{choices}
\end{question}
\answer{5}
